\chapter{Appendix: deployment scenarios}
\label{ch:deployments}

\begin{aside}
The realities of shipping. SSH, kiosk, embedded, container,
WSL --- each environment has quirks.
\end{aside}

\section{SSH}

The remote terminal is the user's local terminal. Capabilities
match the local terminal, mediated by SSH (transparent for
most things; some private modes drop).

Latency is the user's link RTT. Aim for small frames.

\section{Tmux / screen}

The multiplexer is the layer between you and the user's
terminal. CSI-u, OSC-8, image protocols may be filtered.
Document required tmux config (e.g., \code{set -g
allow-passthrough on}).

\section{Kiosk}

A dedicated machine running one app. Full control over
the terminal. Use the best capabilities; assume your
preferred terminal.

\section{Embedded / set-top}

Resource-constrained. Static binary, minimal deps.
Compile for the target ISA. Test on the actual hardware.

\section{Container}

The terminal is whatever the user attached. Default to
moderate caps; let env vars override.

\section{WSL}

Windows Subsystem for Linux. Modern Windows Terminal
supports most features. Old conhost.exe does not.

\section{Cron / non-interactive}

No terminal. \maya{} detects via \code{isatty}; falls
back to no-op output. Good for tools that double as TUI
and pipe.

\section{Recap}

\begin{recap}
\begin{itemize}[leftmargin=*]
  \item SSH: small frames; trust the link.
  \item tmux: passthrough config matters.
  \item Kiosk: assume the best terminal.
  \item Embedded: smallest dep set.
  \item Container, WSL, cron: detect and adapt.
\end{itemize}
\end{recap}

\section*{Workshop}

\begin{enumerate}[leftmargin=*]
  \item Test your app over SSH to a remote box.
  \item Run inside tmux; identify what breaks.
  \item Run with stdout piped to a file.
\end{enumerate}
